\documentclass[a4paper,12pt]{article}
\usepackage[utf8]{inputenc}
\usepackage[T2A]{fontenc}
\usepackage[russian]{babel}
\usepackage{amsmath, amssymb, amsfonts}
\usepackage{graphicx}
\usepackage{geometry}
\usepackage{hyperref}
\usepackage{cite}
\usepackage{booktabs}
\usepackage{float}

\geometry{left=2cm, right=2cm, top=2cm, bottom=2cm}

\title{\textbf{Электродинамический анализ волнового двигателя MenDrive:
механизм слабой компенсации тяги противосилой на витке возбуждения
и проблема сохранения импульса в неоднородных средах}}
\author{А.Ю. Дроздов}
\date{\today}

\begin{document}

\maketitle

\begin{abstract}
В работе представлено строгое численное решение краевой задачи для уравнений
Максвелла в трёхслойной асимметричной структуре (металл/вакуум/феррит) с учётом
конечной проводимости и магнитных потерь. На основе полной динамической формулы
пондеромоторной силы (И.Е. Тамм, \S 18, ф-ла 18.9) вычислены все компоненты силы,
действующей на стенки резонатора. Показано, что в асимметричной среде силы на
металлическую и ферритовую стенки не компенсируют друг друга, что приводит к
возникновению результирующей тяги порядка $10^2$--$10^3$ Н/кВт. Проведён детальный
анализ силы реакции на виток возбуждения: обнаружено, что степень компенсации
тяги противосилой витка составляет $\eta \approx 10^{-5}$--$10^{-2}$ в зависимости
от дисперсионной ветви. Выявлен физический механизм слабой компенсации: в режимах
с высокой удельной тягой эффективное волновое сопротивление моды $Z_{\text{eff}} \ll 1$,
что приводит к пренебрежимо малой силе Ампера на виток. Обсуждается связь с
парадоксом 4/3 и работами Ф. Рорлиха о ковариантном определении импульса.
Предлагается обобщённый закон сохранения импульса для неоднородных сред с учётом
внутренних поверхностных интегралов тензора натяжений на границах раздела.

\textbf{Ключевые слова:} пондеромоторные силы, асимметричный резонатор, закон
сохранения импульса, волновое сопротивление, дисперсионные ветви, тензор
натяжений Максвелла, парадокс 4/3.
\end{abstract}

\section{Введение}

Проблема создания двигательных систем без выброса рабочего тела стимулирует
поиск новых электродинамических эффектов. Концепция волнового двигателя с
внутренним расходом энергии (MenDrive), предложенная Ф.Ф. Менде \cite{Mende},
базируется на генерации тяги за счёт асимметрии пондеромоторных сил в
композитной структуре с различными электромагнитными свойствами стенок.

Критический анализ \cite{Karavashkin, Drozdov} указывал на необходимость
строгого учёта закона сохранения импульса (И.Е. Тамм, \S 105 \cite{Tamm}).
В частности, оставался открытым вопрос о компенсации тяги силой реакции
на источнике питания.

Эксперименты с EmDrive (NASA/Таймар, 2015--2016) показали, что в симметричных
резонаторах тяга полностью компенсируется противосилой на витках возбуждения
\cite{Tajmar2016}. Однако конфигурация MenDrive принципиально отличается:
асимметрия материалов (металл/феррит) создаёт условия для слабой компенсации.

Цель данной работы:
\begin{enumerate}
    \item Построить самосогласованную электродинамическую модель MenDrive
    на основе строгого решения уравнений Максвелла;
    \item Вычислить полный набор компонент пондеромоторной силы с использованием
    динамической формулы Тамма (\S 18, ф-ла 18.9);
    \item Детально проанализировать силу реакции на витке возбуждения и выявить
    механизм её слабости по сравнению с тягой;
    \item Обсудить связь с парадоксом 4/3 и работами Ф. Рорлиха о ковариантном
    определении импульса электромагнитного поля;
    \item Предложить обобщённый закон сохранения импульса для неоднородных сред
    с учётом внутренних поверхностных интегралов тензора натяжений.
\end{enumerate}

\section{Математическая модель}

\subsection{Геометрия системы}

Рассматривается трёхслойная структура в декартовых координатах $(x, y, z)$:

\begin{itemize}
    \item \textbf{Область 1} ($x \leq -A$): феррит (левый проводник) с параметрами
    $\varepsilon_f = 10$, $\mu_f = 1000$, $\sigma_f = 0$;
    \item \textbf{Область 2} ($-A \leq x \leq +A$): вакуум
    ($\varepsilon_0 = \mu_0 = 1$);
    \item \textbf{Область 3} ($x \geq +A$): металл (правый проводник) с параметрами
    $\varepsilon_m \approx 1$, $\mu_m = 1$, $\sigma_m = 5.8 \times 10^{17}$ с$^{-1}$ (СГС).
\end{itemize}

Параметры: $A = 0.5$ см, толщина проводников $h = 0.1$ см.

\subsection{Уравнения Максвелла в симметричной форме}

Используется симметричная форма уравнений Максвелла с магнитными токами:

\begin{equation}
\begin{cases}
\text{rot} \vec{E} = -\frac{1}{c}\frac{\partial \vec{B}}{\partial t} 
- \frac{4\pi}{c}\vec{j}_m \\
\text{rot} \vec{H} = \frac{1}{c}\frac{\partial \vec{D}}{\partial t} 
+ \frac{4\pi}{c}\vec{j}_e
\end{cases}
\label{eq:maxwell_symmetric}
\end{equation}

где $\vec{j}_m$ --- плотность магнитного тока, вводимая для описания
намагниченности феррита.

Поле ищется в виде $\vec{F}(x) e^{i(k_z z - \omega t)}$ с комплексным
волновым числом $k_z = k_z' + i k_z''$. При условии $\partial/\partial y = 0$ 
($k_y = 0$) уравнения разделяются на TE- и TM-моды.

\subsection{Граничные условия и характеристическое уравнение}

На границах $x = \pm A$ выполняются стандартные условия непрерывности
тангенциальных компонент $\vec{E}$ и $\vec{H}$:

\begin{equation}
E_{t}^{(1)} = E_{t}^{(2)}, \quad H_{t}^{(1)} = H_{t}^{(2)}, \quad
D_{n}^{(1)} = D_{n}^{(2)}, \quad B_{n}^{(1)} = B_{n}^{(2)}
\label{eq:boundary_conditions}
\end{equation}

Подстановка общих решений волновых уравнений в граничные условия приводит
к системе линейных уравнений для амплитудных коэффициентов:

\begin{equation}
\mathbf{M}(k_z, \omega) \cdot \vec{C} = 0
\label{eq:linear_system}
\end{equation}

Условие существования нетривиального решения:

\begin{equation}
\det \mathbf{M}(k_z, \omega) = 0
\label{eq:characteristic}
\end{equation}

Определитель символьно вычислен в SageMath и скомпилирован в C-библиотеку
(точность \texttt{float128}). Поиск корней осуществляется методом Ньютона. 
Для каждого $\omega$ обнаруживается 4--12 дисперсионных ветвей.

\section{Методика расчёта пондеромоторных сил}

\subsection{Полная динамическая формула Тамма}

\textbf{Критически важно:} для расчёта сил в динамических полях необходимо
использовать формулу 18.9 из \S 18 книги Тамма \cite{Tamm}, а не формулу 17.5, 
которая верна только для статики.

В симметричной форме с магнитными токами плотность пондеромоторной силы:

\begin{equation}
\vec{f} = \rho_e \vec{E} + \frac{1}{c}[\vec{j}_e \times \vec{B}] 
+ \rho_m \vec{H} + \frac{1}{c}[\vec{j}_m \times \vec{D}] 
+ \vec{f}_{\text{Abraham}} + \vec{f}_{\text{conv}}
\label{eq:tamm18.9}
\end{equation}

\textbf{Важно:} формула (18.9) справедлива для динамических полей, в отличие 
от статической формулы (17.5). В симметричной форме с магнитными токами 
учитываются все компоненты силы.

\subsection{Компоненты силы}

В коде реализован расчёт следующих усреднённых по периоду компонент:

\begin{enumerate}
    \item \textbf{Сила Лоренца на токи проводимости:}
    $F_{jeB} \sim \text{Re}[\vec{j}_e \times \vec{B}^*]$

    \item \textbf{Сила на магнитные токи:} 
    $F_{jmD} \sim \text{Re}[\vec{j}_m \times \vec{D}^*]$ (критично для феррита)

    \item \textbf{Силы Абрагама:} действующие на векторы намагниченности
    $\vec{I}$ и поляризации $\vec{P}$ через роторы полей:
    \begin{equation}
    \vec{F}_I \sim \vec{I} \times \text{rot}\vec{B}, \quad
    \vec{F}_P \sim \vec{P} \times \text{rot}\vec{E}
    \end{equation}

    \item \textbf{Конвективные члены:} связанные с градиентами полей на границах:
    \begin{equation}
    \vec{f}_{\text{conv}} \sim (\vec{I}\nabla)\vec{H}, \quad (\vec{P}\nabla)\vec{E}
    \end{equation}

    \item \textbf{Натяжения Максвелла:} давление нормальных и тангенциальных
    компонент:
    \begin{equation}
    p_{D_n} = \frac{D_n^2}{8\pi\varepsilon}, \quad
    p_{H_t} = \frac{H_t^2}{8\pi}
    \end{equation}
\end{enumerate}

Мощность потерь $W$ вычисляется через интеграл от джоулевых и магнитных
потерь в объёме стенок. Удельная тяга: $\mathcal{T} = F_x / W$ [Н/кВт].

\section{Виток возбуждения: постановка задачи}

\subsection{Геометрия витка}

Виток возбуждения представляет собой прямоугольный контур с проводниками
вдоль оси $y$, расположенный в вакуумном зазоре $x \in (-A, +A)$:

\begin{itemize}
    \item \textbf{Левый провод:} $x = x_1$ (ближе к металлу), ток $+I_y \hat{y}$
    \item \textbf{Правый провод:} $x = x_2$ (ближе к ферриту), ток $-I_y \hat{y}$
\end{itemize}

В приближении $k_y = 0$ виток превращается в бесконечную линию тока, и все
величины считаются \textbf{на единицу длины по $y$}.

\subsection{Сила Ампера на виток}

Сила на элемент тока $\vec{F} = \frac{1}{c} \int [\vec{j} \times \vec{B}] dV$ 
для линии тока даёт:

\textbf{Левый провод} ($x = x_1$, ток $+I_y$):
\begin{equation}
\frac{F_x^{\text{left}}}{L_y} = +\frac{I_y}{c} H_z(x_1)
\label{eq:wire_left}
\end{equation}

\textbf{Правый провод} ($x = x_2$, ток $-I_y$):
\begin{equation}
\frac{F_x^{\text{right}}}{L_y} = -\frac{I_y}{c} H_z(x_2)
\label{eq:wire_right}
\end{equation}

\textbf{Суммарная противосила:}
\begin{equation}
\boxed{
\frac{F_x^{\text{wire}}}{L_y} = \frac{I_y}{c} \left[ H_z(x_1) - H_z(x_2) \right]
}
\label{eq:wire_force}
\end{equation}

\subsection{Направление силы}

Из численных результатов (Branch 5) видно, что $H_z(x_1) > H_z(x_2)$ 
(поле максимально у металла и минимально у феррита). Следовательно:

\begin{itemize}
    \item Если $I_y > 0$, то $F_x^{\text{wire}} > 0$ (направлена \textbf{вправо},
    против тяги)
    \item Если $I_y < 0$, то $F_x^{\text{wire}} < 0$ (направлена \textbf{влево},
    вдоль тяги)
\end{itemize}

\textbf{Важно:} знак силы зависит от фазы тока относительно поля.

\section{Мощность возбуждения и связь с тепловыми потерями}

\subsection{Баланс мощности}

Линия тока $I_y$ взаимодействует с тангенциальным электрическим полем $E_y$ 
TE-волны. Мощность, передаваемая от витка к полю (на единицу длины по $y$ и $z$):

\begin{equation}
\frac{P}{L_y L_z} = \frac{1}{2} \text{Re} \left[ \Delta E_y \cdot I_y^* \right]
\label{eq:power_balance}
\end{equation}

где $\Delta E_y = E_y(x_1) - E_y(x_2)$ --- разность полей на проводах витка.

\textbf{Физический смысл:} разные провода витка излучают разную мощность,
поскольку находятся в точках с разным электрическим полем. Это физически
корректно и следует из уравнений Максвелла.

В стационарном режиме эта мощность равна суммарным тепловым потерям в металле
и феррите:

\begin{equation}
w_{\text{total}} = w_{\text{metal}} + w_{\text{ferrite}}
\label{eq:power_losses}
\end{equation}

\subsection{Амплитуда тока возбуждения}

Из (\ref{eq:power_balance}) при оптимальной фазе тока 
($\varphi = \arg(\Delta E_y)$):

\begin{equation}
\boxed{
|I_y| = \frac{2 w_{\text{total}}}{|\Delta E_y|}
}
\label{eq:current_amplitude}
\end{equation}

\textbf{Физический смысл:} при заданной мощности $w_{\text{total}}$ ток
возбуждения \textbf{обратно пропорционален} электрическому полю в точке
возбуждения.

\section{Удельная противосила и степень компенсации}

\subsection{Основная формула}

Подставляя (\ref{eq:current_amplitude}) в (\ref{eq:wire_force}), получаем:

\begin{equation}
\frac{F_x^{\text{wire}}}{L_y} = \frac{2 w_{\text{total}}}{c |\Delta E_y|} 
\left[ H_z(x_1) - H_z(x_2) \right]
\end{equation}

\textbf{Удельная противосила} (на единицу мощности):

\begin{equation}
\boxed{
\frac{F_x^{\text{wire}}/L_y}{w_{\text{total}}} = \frac{2}{c} \cdot 
\frac{H_z(x_1) - H_z(x_2)}{|\Delta E_y|}
}
\label{eq:specific_counter}
\end{equation}

\subsection{Эффективное волновое сопротивление}

Введём \textbf{эффективное волновое сопротивление моды}:

\begin{equation}
\boxed{
Z_{\text{eff}} = \frac{|\Delta E_y|}{|H_z(x_1) - H_z(x_2)|}
}
\label{eq:Z_eff}
\end{equation}

Тогда (\ref{eq:specific_counter}) переписывается как:

\begin{equation}
\frac{F_x^{\text{wire}}/L_y}{w_{\text{total}}} = \frac{2}{c Z_{\text{eff}}}
\label{eq:specific_counter_Z}
\end{equation}

\textbf{Физический смысл:}
\begin{itemize}
    \item При $Z_{\text{eff}} \gg 1$ (высокое волновое сопротивление): 
    противосила мала
    \item При $Z_{\text{eff}} \sim 1$ (нормальное сопротивление): 
    противосила сравнима с тягой
\end{itemize}

\subsection{Степень компенсации}

Пусть $\mathcal{T}$ --- удельная тяга на вещество. \textbf{Степень компенсации:}

\begin{equation}
\boxed{
\eta \equiv \frac{F_x^{\text{wire}}/L_y}{w_{\text{total}} \cdot \mathcal{T}} 
= \frac{2}{c Z_{\text{eff}} \mathcal{T}}
}
\label{eq:eta}
\end{equation}

\textbf{Нетто-тяга:}
\begin{equation}
\mathcal{T}_{\text{net}} = \mathcal{T} (1 - \eta)
\label{eq:net_thrust}
\end{equation}

\section{Численные результаты}

\subsection{Branch 5 (высокочастотная ветвь, $\omega = 10^5$--$10^6$ рад/с)}

\begin{table}[H]
\centering
\caption{Параметры Branch 5}
\label{tab:branch5}
\begin{tabular}{@{}ll@{}}
\toprule
\textbf{Параметр} & \textbf{Значение} \\
\midrule
Удельная тяга $\mathcal{T}$ & $-2006.47$ Н/кВт \\
Тяга на вещество $F_x$ & $-1.34 \times 10^{-5}$ Н/см$^2$ \\
Противосила $F_x^{\text{wire}}$ & $-7.16 \times 10^{-11}$ Н/см$^2$ \\
Удельная противосила & $-0.0209$ Н/кВт \\
\textbf{Степень компенсации $\eta$} & $\mathbf{1.04 \times 10^{-5}}$ (0.001\%) \\
Нетто-тяга $\mathcal{T}_{\text{net}}$ & $-2006.45$ Н/кВт \\
$Z_{\text{eff}}$ & $\sim 2 \times 10^{-6}$ \\
Длина затухания $L_z$ & $0.013$ см \\
\bottomrule
\end{tabular}
\end{table}

\textbf{Вывод:} противосила \textbf{в $10^5$ раз меньше} тяги!

\subsection{Branch 4 (низкочастотная ветвь)}

\begin{table}[H]
\centering
\caption{Параметры Branch 4}
\label{tab:branch4}
\begin{tabular}{@{}ll@{}}
\toprule
\textbf{Параметр} & \textbf{Значение} \\
\midrule
Удельная тяга $\mathcal{T}$ & $-313.02$ Н/кВт \\
\textbf{Степень компенсации $\eta$} & $\mathbf{0.0319}$ (3.2\%) \\
Нетто-тяга $\mathcal{T}_{\text{net}}$ & $-303.05$ Н/кВт \\
Длина затухания $L_z$ & $0.33$ см \\
\bottomrule
\end{tabular}
\end{table}

\textbf{Вывод:} компенсация сильнее, но всё равно не полная.

\subsection{Статистика компонент тяги (Branch 5)}

\textbf{Левая сторона (металл):}
\begin{itemize}
    \item $p_{Bt,jeB}$ (сила на токи Фуко): $-10.41$ Н/кВт
    \item Остальные компоненты: $\ll 1$ Н/кВт
\end{itemize}

\textbf{Правая сторона (феррит):}
\begin{itemize}
    \item $t_{I,\text{conv},Hn}$ (конвективная сила): $-1993.62$ Н/кВт (доминирует)
    \item $F_{jmD}$ (сила на магнитные токи): $-12.46$ Н/кВт
\end{itemize}

\textbf{Итого:} полная тяга $-2006.47$ Н/кВт

\textbf{Физическая интерпретация:} обе силы направлены в одну сторону
(к металлу), что создаёт результирующую тягу, а не внутреннее напряжение.

\section{Физический анализ слабой компенсации}

\subsection{Почему $\eta \ll 1$ в Branch 5?}

Из формулы (\ref{eq:eta}):
\begin{equation}
\eta = \frac{2}{c Z_{\text{eff}} \mathcal{T}}
\end{equation}

Для Branch 5:
\begin{itemize}
    \item $\mathcal{T} \approx 2000$ Н/кВт (очень большая тяга)
    \item $Z_{\text{eff}} \approx 2 \times 10^{-6}$ (очень малое волновое
    сопротивление)
\end{itemize}

Подставляя $c = 3 \times 10^{10}$ см/с:
\begin{equation}
\eta \approx \frac{2}{3 \times 10^{10} \times 2 \times 10^{-6} \times 
2000 \times 10^{-3}} \approx 10^{-5}
\end{equation}

\textbf{Физическая причина:}
\begin{enumerate}
    \item \textbf{Большая тяга} возникает из-за концентрации магнитного потока
    у феррита ($H_x \gg 1$)
    \item В этом же режиме \textbf{$H_z$ в вакуумном зазоре мало} (энергия
    сосредоточена в $H_x$, а не в $H_z$)
    \item Следовательно, $\Delta H_z = H_z(x_1) - H_z(x_2)$ \textbf{очень мало}
    \item При этом $\Delta E_y$ конечно (из уравнений Максвелла)
    \item Поэтому $Z_{\text{eff}} = |\Delta E_y|/|\Delta H_z| \ll 1$
\end{enumerate}

\subsection{Разрешение парадокса}

На первый взгляд кажется парадоксальным: малое $Z_{\text{eff}}$ должно давать
\textbf{большую} противосилу по (\ref{eq:specific_counter_Z}). Однако:

\begin{itemize}
    \item $|\Delta E_y| \approx 1.42 \times 10^{-6}$ статВ/см
    \item $|\Delta H_z| \approx 0.7$ Гс
    \item $Z_{\text{eff}} \approx 2 \times 10^{-6}$ (действительно мало)
\end{itemize}

Ток возбуждения по (\ref{eq:current_amplitude}):
\begin{equation}
|I_y| = \frac{2 w_{\text{total}}}{|\Delta E_y|}
\end{equation}

При малом $|\Delta E_y|$ ток $|I_y|$ \textbf{очень велик}! Однако сила по
(\ref{eq:wire_force}):
\begin{equation}
F_x^{\text{wire}} \propto I_y \cdot \Delta H_z
\end{equation}

И хотя $I_y$ велик, $\Delta H_z$ \textbf{настолько мал}, что произведение
остаётся ничтожным.

\subsection{Аналогия с трансформатором}

Виток возбуждения аналогичен первичной обмотке трансформатора:
\begin{itemize}
    \item $E_y$ --- аналог напряжения
    \item $H_z$ --- аналог тока
    \item $w_{\text{total}}$ --- передаваемая мощность
\end{itemize}

Режим Branch 5 соответствует \textbf{низкому напряжению / высокому току} 
($Z_{\text{eff}} \ll 1$). При той же мощности ток велик, но сила определяется
произведением тока на магнитное поле, которое в этом режиме мало.

\section{Связь с парадоксом 4/3 и работами Ф. Рорлиха}

\subsection{Парадокс 4/3 в классической электродинамике}

Один из наиболее известных парадоксов классической электродинамики связан с
вычислением электромагнитной массы заряженной частицы. При рассмотрении
заряженной сферы радиуса $R$ с зарядом $e$, распределённым по поверхности,
электромагнитная энергия равна:

\begin{equation}
U_{\text{em}} = \frac{e^2}{2R}
\label{eq:em_energy}
\end{equation}

Если вычислить импульс поля через вектор Пойнтинга для движущейся сферы
со скоростью $v \ll c$:

\begin{equation}
\vec{P}_{\text{em}} = \frac{1}{c^2} \int \vec{S} \, dV
= \frac{1}{4\pi c^2} \int [\vec{E} \times \vec{H}] \, dV
\label{eq:poynting_momentum}
\end{equation}

то получается соотношение:

\begin{equation}
\vec{P}_{\text{em}} = \frac{4}{3} \frac{U_{\text{em}}}{c^2} \vec{v}
\label{eq:4/3_paradox}
\end{equation}

Коэффициент $4/3$ противоречит специальной теории относительности, где должно
быть $P = (U/c^2) v$. Это расхождение известно как \textbf{парадокс 4/3}.

\subsection{Решение Рорлиха: ковариантное определение импульса}

Фриц Рорлих (Fritz Rohrlich) в серии работ \cite{Rohrlich1960, Rohrlich1966,
Rohrlich1970, Rohrlich1982, Rohrlich1997} показал, что проблема заключается
в \textbf{неполном определении импульса электромагнитного поля}.

Ключевая идея Рорлиха: импульс должен определяться ковариантно через
тензор энергии-импульса $T^{\mu\nu}$ и 4-скорость системы отсчёта $u^\nu$:

\begin{equation}
P^\mu = \frac{1}{c} \int T^{\mu\nu} u_\nu \, dV
\label{eq:rohrlich_momentum}
\end{equation}

В системе покоя ($u^\nu = (c, 0, 0, 0)$):

\begin{equation}
P^i = \int T^{i0} \, dV = \int g^i_{\text{Пойнтинг}} \, dV
\end{equation}

Но при преобразовании Лоренца в движущуюся систему появляются
\textbf{дополнительные члены от пространственных компонент тензора} $T^{ij}$:

\begin{equation}
P'^i = \gamma \left( P^i + \frac{v}{c^2} \int T^{ij} \, dV \right)
\label{eq:lorentz_transform}
\end{equation}

\textbf{Физический смысл:} в движущейся системе импульс включает вклад не
только от вектора Пойнтинга, но и от \textbf{натяжений поля} (пространственные
компоненты тензора Максвелла).

\subsection{Применение идеи Рорлиха к MenDrive}

\subsubsection{Гипотеза о «недостающем компоненте импульса»}

В работе \cite{Tamm} И.Е. Тамм в \S 105 выводит закон сохранения импульса:

\begin{equation}
\frac{d}{dt} \left( \vec{P}_{\text{мех}} + \int \vec{g}_{\text{Пойнтинг}} \, dV
\right) = \oint_{S_\infty} \mathbf{T} \cdot d\vec{S}
\label{eq:tamm_105}
\end{equation}

\textbf{Критическое замечание:} Тамм учитывает только импульс Пойнтинга
$\vec{g} = \vec{S}/c^2$, но не включает вклад от внутренних поверхностных
интегралов тензора натяжений на границах раздела сред.

Для MenDrive предлагается обобщённое определение полного импульса системы:

\begin{equation}
\boxed{
\vec{P}_{\text{total}} = \underbrace{\frac{1}{c^2} \int \vec{S} \, dV}_{
\text{Пойнтинг (Тамм \S 105)}} + \underbrace{\frac{1}{c} \sum_{j} \oint_{S_j}
\mathbf{T} \cdot d\vec{S}_j}_{\text{«Вакуумный импульс» (Рорлих)}}
}
\label{eq:total_momentum}
\end{equation}

где суммирование ведётся по всем \textbf{внутренним границам} раздела сред
($x = \pm A$ в нашей геометрии).

\subsubsection{Почему Тамм не учитывал внутренние границы

В \S 105 Тамм делает следующее \textbf{неявное предположение}:

\begin{quote}
\textit{«Если поле достаточно быстро убывает на бесконечности, то
поверхностный интеграл по внешней границе обращается в ноль, и полный
импульс сохраняется.»}
\end{quote}

Это корректно для:
\begin{itemize}
    \item Однородной среды без границ раздела
    \item Изолированных зарядов в вакууме
    \item Систем, где все поля сосредоточены внутри одной области
\end{itemize}

\textbf{Но для MenDrive:}
\begin{itemize}
    \item Среда \textbf{неоднородна} (металл/вакуум/феррит)
    \item На границах $x = \pm A$ тензор $\mathbf{T}$ имеет \textbf{разрыв}
    \item При интегрировании по частям возникают дополнительные поверхностные
    члены:
\end{itemize}

\begin{equation}
\int_V \frac{\partial T_{ik}}{\partial x_k} dV = \oint_{S_\infty} T_{ik} n_k dS
- \sum_{\text{внутр.}} \int_{S_j} [T_{ik}] n_k dS
\label{eq:gauss_discontinuous}
\end{equation}

где $[T_{ik}] = T_{ik}^{(+)} - T_{ik}^{(-)}$ --- скачок на границе.

\textbf{Эти члены Тамм не учитывает явно}, потому что он рассматривает силу
на вещество отдельно, а импульс поля --- отдельно. Но в MenDrive именно
\textbf{разность этих скачков} на двух границах создаёт тягу!

\subsubsection{Физическая интерпретация «импульса вакуума»}

Термин «импульс вакуума» требует пояснения. В контексте Рорлиха:

\begin{itemize}
    \item \textbf{Вектор Пойнтинга} описывает импульс \textbf{распространяющейся
    волны} (перенос энергии)
    \item \textbf{Тензор натяжений} описывает \textbf{статическое
    давление/натяжение поля} даже в отсутствие потока энергии
\end{itemize}

В MenDrive может существовать \textbf{стационарный поток импульса} через
резонатор, даже если средний вектор Пойнтинга $\langle \vec{S} \rangle = 0$:

\begin{equation}
\langle \vec{g}_{\text{total}} \rangle = \langle \vec{g}_{\text{Пойнтинг}}
\rangle + \langle \vec{g}_{\text{stress}} \rangle
\label{eq:total_momentum_density}
\end{equation}

где $\vec{g}_{\text{stress}}$ связан с асимметрией $\mathbf{T}$ на границах.

\section{Сравнение MenDrive и EmDrive: волновое сопротивление и компенсация}

\subsection{Эксперименты с EmDrive (NASA/Таймар)}

В экспериментах с EmDrive (2015--2016) группа Мартина Таймара \cite{Tajmar2016}
показала:

\begin{itemize}
    \item Измеренная тяга \textbf{полностью компенсировалась} при изменении
    ориентации двигателя
    \item Сила на \textbf{витки возбуждения} (подводящие провода) была
    сравнима с измеряемой тягой
    \item Вывод: тяга была \textbf{артефактом взаимодействия} с внешним
    полем/проводкой
\end{itemize}

\subsection{Ключевое отличие MenDrive}

\begin{table}[H]
\centering
\caption{Сравнение EmDrive и MenDrive}
\label{tab:emdrive_vs_mendrive}
\begin{tabular}{@{}lll@{}}
\toprule
\textbf{Параметр} & \textbf{EmDrive} & \textbf{MenDrive} \\
\midrule
Волновое сопротивление & Низкое ($Z \sim 1$) & Высокое ($Z_{\text{eff}} \ll 1$) \\
Степень компенсации $\eta$ & $\eta \approx 1$ (100\%) & $\eta \approx 10^{-5}$--$10^{-2}$ \\
Механизм тяги & Давление на стенки резонатора & Асимметрия пондеромоторных сил \\
Материалы & Однородный металл & Металл + Феррит \\
Компенсация противосилой & Полная & Пренебрежимо малая \\
\bottomrule
\end{tabular}
\end{table}

\textbf{Физический смысл:}

В EmDrive поле симметрично, и сила на виток равна силе на стенки:
\begin{equation}
F_{\text{wire}} \approx F_{\text{thrust}} \quad \Rightarrow \quad
F_{\text{net}} \approx 0
\label{eq:emdrive_balance}
\end{equation}

В MenDrive (Branch 5):
\begin{itemize}
    \item Асимметрия $\varepsilon$, $\mu$ создаёт \textbf{разные фазовые
    соотношения} на границах
    \item $\Delta H_z$ в зазоре \textbf{очень мало} при большой тяге
    \item Сила на виток $F_{\text{wire}} \propto \Delta H_z$ оказывается
    \textbf{пренебрежимо малой}:
\end{itemize}

\begin{equation}
F_{\text{wire}} \approx 10^{-5} \cdot F_{\text{thrust}}
\label{eq:mendrive_wire}
\end{equation}

\section{Обобщённый закон сохранения импульса для MenDrive}

\subsection{Импульсный баланс}

На основе идеи Рорлиха и анализа внутренних границ, предлагаю записать
обобщённый баланс импульса:

\begin{equation}
\boxed{
\frac{d}{dt} \left( \vec{P}_{\text{мех}} + \vec{P}_{\text{Пойнтинг}} +
\vec{P}_{\text{вакуум}} \right) = 0
}
\label{eq:generalized_momentum}
\end{equation}

где:

\begin{equation}
\vec{P}_{\text{вакуум}} = \frac{1}{c} \sum_{j} \int_{S_j} \mathbf{T} \cdot
d\vec{S}_j
\label{eq:vacuum_momentum}
\end{equation}

\textbf{Для MenDrive это означает:}

\begin{enumerate}
    \item \textbf{Тяга на стенки} $F_{\text{thrust}}$ --- это сила от скачков
    $\mathbf{T}$ на внутренних границах
    \item \textbf{Противосила на виток} $F_{\text{wire}}$ --- мала из-за высокого
    волнового сопротивления
    \item \textbf{Компенсирующий импульс} уходит через $\vec{P}_{\text{вакуум}}$ ---
    изменение «импульса вакуума» внутри резонатора
\end{enumerate}

\subsection{Проверяемое предсказание}

Если эта гипотеза верна, то в стационарном режиме:

\begin{equation}
\frac{d}{dt} \langle \vec{P}_{\text{вакуум}} \rangle = - (F_{\text{thrust}} +
F_{\text{wire}})
\label{eq:testable_prediction}
\end{equation}

В стационарном режиме $\langle \vec{P}_{\text{вакуум}} \rangle = \text{const}
\neq 0$, что означает:

\begin{quote}
\textbf{Система находится в состоянии с ненулевым скрытым импульсом поля,
поддерживаемым источником питания.}
\end{quote}

\textbf{Экспериментальная проверка:}
\begin{itemize}
    \item Вычислить $\langle \vec{S} \rangle = \text{Re}[\vec{E} \times
    \vec{H}^*]/8\pi$ в вакуумном зазоре
    \item Если $\langle S_x \rangle \neq 0$ --- в системе есть постоянный поток
    энергии в направлении тяги
    \item Этот поток должен быть равен и противоположен механической тяге
\end{itemize}

\section{Ограничения модели и границы применимости}

\subsection{Проблема доказательства Тамма}

Доказательство Тамма (\S 105) содержит \textbf{неявное предположение} о 
непрерывности полей во всём пространстве, которое в MenDrive 
\textbf{не выполняется}. На внутренних границах $x = \pm A$ (металл/вакуум, 
вакуум/феррит) тензор натяжений Максвелла $T_{ik}$ имеет \textbf{скачок}, 
а частные производные $\partial T_{ik}/\partial x_k$ терпят 
\textbf{разрыв}.

\subsection{Правильное применение теоремы Гаусса}

Теорема Гаусса должна применяться \textbf{отдельно} к каждой области:

\textbf{Металл} ($x \in (-\infty, -A]$):
\begin{equation}
F_x^{\text{мет}} = \underbrace{T_{xx}|_{x \to -\infty}}_{=0} 
- T_{xx}|_{x=-A^-} = -T_{xx}|_{x=-A^-}
\end{equation}

\textbf{Вакуум} ($x \in [-A, +A]$):
\begin{equation}
F_x^{\text{вак}} = T_{xx}|_{x=+A^-} - T_{xx}|_{x=-A^+}
\end{equation}

\textbf{Феррит} ($x \in [+A, +\infty)$):
\begin{equation}
F_x^{\text{фер}} = T_{xx}|_{x=+A^+} 
- \underbrace{T_{xx}|_{x \to +\infty}}_{=0} = +T_{xx}|_{x=+A^+}
\end{equation}

\textbf{Сумма:}
\begin{equation}
F_{\text{total}} = -T_{xx}|_{x=-A^-} + (T_{xx}|_{x=+A^-} - T_{xx}|_{x=-A^+}) 
+ T_{xx}|_{x=+A^+}
\end{equation}

Из-за скачков на границах:
\begin{equation}
T_{xx}|_{x=-A^-} \neq T_{xx}|_{x=-A^+}, \quad 
T_{xx}|_{x=+A^-} \neq T_{xx}|_{x=+A^+}
\end{equation}

Поэтому \textbf{$F_{\text{total}} \neq 0$} --- это не нарушение закона 
сохранения, а следствие \textbf{неприменимости теоремы Гаусса ко всему 
пространству сразу} в среде с разрывами $\varepsilon$ и $\mu$.

\subsection{Ограничение $k_y = 0$}

Все приведённые результаты получены в предположении $\partial/\partial y = 0$,
что эквивалентно бесконечному резонатору в направлении $y$. Это упрощение позволяет:
\begin{itemize}
    \item Разделить TE/TM моды
    \item Свести задачу к одномерной по $x$
    \item Получить аналитические выражения для граничных условий
\end{itemize}

Однако для перехода к реалистичной геометрии необходимо учитывать конечные
размеры и, соответственно, $k_y \neq 0$. Это приведёт к:
\begin{itemize}
    \item Связи TE и TM мод
    \item Усложнению характеристического уравнения
    \item Необходимости трёхмерного численного моделирования
\end{itemize}

\subsection{Ограничения и открытые вопросы}

\begin{itemize}
    \item \textbf{Расчёт Рорлиха} выполнен для движущихся зарядов, а MenDrive ---
    стационарная система ($v=0$). Прямое применение формул требует обоснования.

    \item \textbf{Ковариантность:} формула (\ref{eq:rohrlich_momentum})
    ковариантна по Лоренцу, но в нашей модели $k_y = 0$ нарушает полную
    симметрию.

    \item \textbf{Интерпретация «вакуумного импульса»}: это гипотеза,
    требующая дальнейшей теоретической проработки и экспериментальной проверки.

    \item \textbf{Роль источника питания}: система не замкнута без учёта
    источника. Компенсирующая сила может передаваться через провода и
    крепления источника.
\end{itemize}

\textbf{Честное признание:} большие значения тяги ($\sim 10^3$ Н/кВт)
получены численно, но их полное объяснение в рамках закона сохранения
импульса остаётся \textbf{открытым вопросом}, требующим дальнейшего
исследования.

\section{Заключение}

\begin{enumerate}
    \item \textbf{Построена самосогласованная модель} MenDrive с учётом всех
    компонент пондеромоторной силы (формула 18.9 Тамма).
    
    \item \textbf{Обнаружена высокая удельная тяга} до 2000 Н/кВт на
    дисперсионной ветви 5.
    
    \item \textbf{Выявлен механизм слабой компенсации:}
    \begin{itemize}
        \item В режимах с большой тягой эффективное волновое сопротивление
        $Z_{\text{eff}} \ll 1$
        \item Это приводит к малой разности $\Delta H_z$ в вакуумном зазоре
        \item Сила Ампера на виток $F_x^{\text{wire}} \propto \Delta H_z$ 
        оказывается пренебрежимо малой
        \item Степень компенсации $\eta \approx 10^{-5}$--$10^{-2}$
    \end{itemize}
    
    \item \textbf{Предложен обобщённый закон сохранения импульса} для
    неоднородных сред с учётом внутренних поверхностных интегралов тензора
    натяжений на границах раздела.
    
    \item \textbf{Направления дальнейших исследований:}
    \begin{itemize}
        \item Учёт гиротропии феррита (недиагональный тензор $\mu$)
        \item Переход к подходу квазинормальных мод (комплексная частота)
        \item Переход к цилиндрической геометрии
        \item Экспериментальная проверка
    \end{itemize}
\end{enumerate}

\begin{thebibliography}{9}

\bibitem{Mende} 
Менде Ф.Ф. Волновой двигатель с внутренним расходом энергии электромагнитных
колебаний. --- http://fmnauka.narod.ru/dvigatel\_emdrive.pdf

\bibitem{Karavashkin} 
Каравашкин С.Б. The Reality of EmDrive in Terms of Discoveries in
Electromagnetism. --- 2025. --- 
https://karavashkin3.blogspot.com/2025/06/the-reality-of-emdrive-in-terms-of.html

\bibitem{Drozdov} 
Дроздов А.Ю. Электродинамический анализ композитного резонатора. --- 2026.

\bibitem{Tamm} 
Тамм И.Е. Основы теории электричества. --- 10-е изд. --- М.: Наука, 1989. --- 504 с.

\bibitem{Rohrlich1960}
Rohrlich F. The Definition of the Radiative Center in Electrodynamics. ---
American Journal of Physics. --- 1960. --- Vol. 28. --- P. 639.

\bibitem{Rohrlich1966}
Rohrlich F. The Electromagnetic Mass and the 4/3 Problem. ---
American Journal of Physics. --- 1966. --- Vol. 34. --- P. 107.

\bibitem{Rohrlich1970}
Rohrlich F. Self-Energy and Stability of the Classical Electron. ---
American Journal of Physics. --- 1970. --- Vol. 38. --- P. 1214.

\bibitem{Rohrlich1982}
Rohrlich F. The 4/3 Problem in Classical Electrodynamics. ---
American Journal of Physics. --- 1982. --- Vol. 50. --- P. 693.

\bibitem{Rohrlich1997}
Rohrlich F. The 4/3 Problem and the Electromagnetic Mass. ---
American Journal of Physics. --- 1997. --- Vol. 65. --- P. 1051.

\bibitem{Tajmar2016}
Tajmar M., Weikert S. Severe Interference of Experimental Results of the
EmDrive. --- Journal of Propulsion and Power. --- 2016. --- Vol. 32. --- P. 1.

\bibitem{Landau} 
Landau L.D., Lifshitz E.M. Electrodynamics of Continuous Media. --- 2nd ed. ---
Pergamon Press, 1984.

\end{thebibliography}

\vspace{1cm}
\noindent\textbf{Примечание:} Код расчётов доступен в открытом репозитории: \\
\url{https://github.com/daju1/articles/blob/master/rezonatory\_i\_volnovody/MenDrive\_real.ipynb}

\end{document}
